\hypertarget{appendix-341296fb3f64}{}
\section*{Robust geometric predicates}
\chapter{Robust geometric predicates}\label{appendix-341296fb3f64-robust-geometric-predicates}

\textbf{Curriculum depth:} Introduced · \textbf{First seen:} Phase6 · \textbf{Latest development:} Phase6

\section{Definition and mechanism}\label{appendix-341296fb3f64-definition-and-mechanism}

A geometric predicate decides a category from a sign, such as orientation. Near cancellation, ordinary floating arithmetic may not certify that sign. Adaptive arithmetic adds work only when justified error bounds leave uncertainty. Exact stored-input decisions cannot restore rounded-away physical information.

A triangle can have an accurately small area estimate yet an uncertified orientation sign. Declaring every small value zero changes topology instead of proving a decision.

\section{Prerequisites}\label{appendix-341296fb3f64-prerequisites}

\begin{itemize}
\tightlist
\item
  \hyperref[appendix-2391db8052b1]{Cross products and oriented area} --- Oriented area supplies face normals, frame signs and geometric determinants.
\item
  \href{../01-Foundations-I/main.pdf\#nameddest=appendix-5aef4f13483a}{Dot and cross-product rounding (01-Foundations-I, appendix-5aef4f13483a)} --- Cancellation and rounded inputs explain when a determinant sign or finite measurement is uncertain.
\end{itemize}

\section{Used by}\label{appendix-341296fb3f64-used-by}

Later curriculum sources can deepen this same node; no future coverage is marked complete.

\begin{itemize}
\tightlist
\item
  \href{../04-Visual-Computing-II/main.pdf\#nameddest=appendix-a9ba01fa72a6}{Collision barriers and continuous detection (04-Visual-Computing-II, appendix-a9ba01fa72a6)} \#\# Recall and next depth
\end{itemize}

Distinguish tolerance-based geometry, measurements and certified predicates; the Phase 6 lab implements only ordinary tests. Full production proofs, target-specific integration and learner mastery require further evidence.

\section{Source and navigation}\label{appendix-341296fb3f64-source-and-navigation}

\href{Continuous\%20Notes/Computing\%20Foundations\%20-\%20From\%20Bits\%20to\%20C++\%20and\%20Systems\%20APIs\#The\%20sensor\%20needs\%20coordinates\%20inside\%20the\%20triangle}{Teaching narrative} · \href{../01-Foundations-I/main.pdf\#nameddest=appendix-fd090d982a15-phase-6---parts-501-600}{Coverage and source gap (01-Foundations-I, appendix-fd090d982a15-phase-6-\/-\/-parts-501-600)} · \href{../01-Foundations-I/main.pdf\#nameddest=appendix-a103e7846d21-phase-6---source-analysis-and-verification}{Evidence and limits (01-Foundations-I, appendix-a103e7846d21-phase-6-\/-\/-source-analysis-and-verification)} · \hyperref[appendix-4ad47bdd269a]{Dependency Map}

\section{Phase 11 camera development}\label{appendix-341296fb3f64-phase-11-camera-development}

Adaptive orientation and incircle predicates protect discrete sign decisions from floating-point error by refining arithmetic when an error bound cannot certify the result. Camera estimation has a different uncertainty: noisy pixels and weak observability can leave pose uncertain even when every sign is exact. Predicate robustness does not improve measurement accuracy, and a low reprojection residual does not guarantee every combinatorial decision was reliable. Keep layers distinct: use robust predicates for topology, uncertainty weights for image residuals, and Jacobian spectra for weak directions. Link: \href{Continuous\%20Notes/Computing\%20Foundations\%20-\%20From\%20Bits\%20to\%20C++\%20and\%20Systems\%20APIs\#A\%20camera\%20model\%20must\%20close\%20the\%20loop}{A camera model must close the loop}; research: \href{Continuous\%20Notes/Computing\%20Foundations\%20-\%20From\%20Bits\%20to\%20C++\%20and\%20Systems\%20APIs\#Four\%20deeper\%20links\%20from\%20the\%20TUM\%20lecture\%20review}{Four deeper links from the TUM lecture review}.
